\subsection{Standing hypotheses and notation}
\label{sec:ch34-hypotheses}

Throughout the chapter \(k\) is an algebraically closed field equipped with
the trivial absolute value.  Following \cite[\S1]{BJ22}, and departing in
this part from the default of the conventions sheet, a \emph{variety} over
\(k\) means a reduced separated scheme of finite type; it need not be
irreducible.  Thus \(X\) is a projective \(k\)-variety with irreducible
components \((X_\alpha)_{\alpha\in A}\), a finite family.  We write
\(N^1(X)=N^1(X)_{\mathbf R}\) for the real N\'eron--Severi space and
\(\Pic(X)_{\mathbf Q}=\Pic(X)\otimes\mathbf Q\).  Statements requiring a
function field, a normalization, or a birational model assume in addition
that \(X\) is integral, and say so.

The analytic vocabulary is that of \Cref{chap:trivial-analytification}.  We
write \(s(v)\) for the support of \(v\in X^{\An}\)
(\Cref{def:triv-support}), \(c(v)\) for its center, and
\begin{equation}
\label{eq:ch34-loci}
 X^{\triv}\subset X^{\An},\qquad
 X^{\val}\subset X^{\An},\qquad
 X^{\divis}\subset X^{\val},\qquad
 X^{\lin}\subset X^{\val}
\end{equation}
for, respectively: the trivial semivaluations \(v_{Y,\triv}\) of the
irreducible subvarieties \(Y\subset X\); the semivaluations of full support,
that is, those whose support is an irreducible component
\eqref{eq:triv-full-support-locus}; the divisorial points, namely
\(v_{\triv,\alpha}:=v_{X_\alpha,\triv}\) together with the positive rational
multiples \(t\ord_E\) of prime-divisor valuations over the components; and
the valuations of linear growth, studied systematically in
\Cref{chap:bj-linear-growth}.  A point of
\(X^{\An}\setminus X^{\val}\) is said to have \emph{proper support}.
